\chapter{Experimental methodology}
\label{cha:methodology}

The claims of this thesis are comparative, so the grid varies one factor at a time from a fixed baseline to attribute any measured difference to the factor changed. Because neither the system nor the judge measuring it is deterministic, every cell is repeated on two independent axes (\Cref{sec:methodology-stability}).

\begin{definition}[Cell]
    \label{def:cell}
    A \emph{cell} of an experimental grid is one combination of language model, system variant and scenario.
\end{definition}

\section{Research questions}
\label{sec:methodology-research-questions}

The hypotheses in \Cref{sec:introduction-goal} translate into three concrete research questions.

\begin{description}
    \item[RQ1] By how much does enabling the deterministic totaller reduce energy and macronutrient error?
    \item[RQ2] By how much does enabling the critique-and-refine loop improve adherence to the stated constraints and preferences, and what effect does it have on energy and macronutrient error, on its own and on top of the totaller?
    \item[RQ3] How does plan quality vary with the choice of language model and with its reasoning-effort level, and can raising the effort level substitute for a stronger model?
\end{description}

RQ1~and RQ2~are answered by a \emph{component ablation} experiment, \textbf{Experiment~1} (\Cref{sec:methodology-variants}), which holds the language model fixed and varies only the totaller and reflection flags. RQ3~is answered by a \emph{model comparison} experiment, \textbf{Experiment~2} (\Cref{sec:methodology-model-grid}), which holds the system configuration fixed at its strongest setting and varies the model and its reasoning-effort level together.

\section{Experiment 1: component ablation}
\label{sec:methodology-variants}

The component ablation experiment uses the four-way grid implied by the two independent flags in \Cref{sec:design-variant-config}, shown in \Cref{fig:diagram-variants}. Food knowledge base access is not a varied flag (\Cref{sec:vision-food-lookup-baseline}): \emph{baseline} here denotes a system that already looks up real foods and merely lacks the two optional mechanisms under study.

Each of the four variants is run against every scenario in \Cref{sec:methodology-scenarios} with one fixed language model, giving 60 runs in total for this experiment. The following parameters are held constant across both experiments: three repetitions per cell, three scoring passes per judge per run, and a critique-and-refine loop budget of three refinement passes (\Cref{def:iterations}).

\section{Experiment 2: model and reasoning-effort grid}
\label{sec:methodology-model-grid}

The model comparison experiment fixes the system at the full variant and varies two factors at once: the base model and the reasoning-effort level it is asked to run at. Four base models are compared, each at three effort levels, giving twelve configurations. Where a provider's default effort already falls at one of the three levels, that default is used, so a configuration is identified by a model together with the effort it actually ran at.

Two properties of this design matter for interpreting the results. Within one base model the three configurations differ only in effort, so a difference between them is attributable to deliberation alone. Across base models the effort labels are each vendor's own, with no guarantee that one vendor's \emph{high} corresponds to another's, so the aggregate comparison across effort levels in \Cref{sec:results-effort} groups vendor-specific labels into three coarse buckets and indicates direction, not a calibrated measurement. Each of the twelve configurations is run against all five scenarios, giving 180 runs. Fifteen coincide with the full-variant cell of the ablation experiment and are stored and scored once, so the two experiments together comprise 225 distinct runs.

The reasoning-effort setting is applied only to the agents whose output affects plan correctness -- the nutrition agent, the critic and the refiner (\Cref{sec:design-agents}). The judges run at their own providers' defaults regardless of which configuration produced the plan, so the measuring instrument does not change with the object of measurement.

\section{Scenarios and profiles}
\label{sec:methodology-scenarios}

Each evaluated run pairs a fixed user profile with a scenario: an optional free-form session request and, where applicable, a set of scenario-specific soft criteria scored by the qualitative judges (\Cref{sec:methodology-qualitative}). The five scenarios in \Cref{tab:methodology-scenarios} were chosen to differ in what they demand of the system, not to sample a population: one asks for nothing beyond a coherent day of food, two set a hard exclusion against an explicit user request for something that would violate it, and two impose interacting chronic conditions whose dietary guidance pulls in several directions at once.

\begin{table}[htbp]
    \centering
    \begin{tabular}{p{0.26\textwidth}p{0.66\textwidth}}
        \toprule
        Scenario & What it demands of the system \\
        \midrule
        regular & Healthy adult, no session request; scored only on general recipe coherence, so it isolates baseline planning quality with nothing else competing for attention. \\
        lactose-intolerant-athlete-wants-cheesecake & Lactose-intolerant athlete explicitly requests cheesecake, testing whether an allergen exclusion survives a direct, conflicting user request, on top of an unusually high energy target. \\
        vegetarian-allergic & Vegetarian with a nut allergy requests peanut-butter-and-jam toast: the same conflicting-request pattern, applied to a diet pattern and a second, independent allergen at the same time. \\
        diabetes-hypertension & Type~2 diabetes and hypertension, no session request; scored against DASH-pattern sodium and potassium targets and diabetes-specific micronutrient considerations that constrain each other. \\
        dyslipidemia-obesity & Dyslipidaemia and obesity, requests pizza; combines the conflicting-request pattern with cardiometabolic micronutrient and macronutrient targets. \\
        \bottomrule
    \end{tabular}
    \caption{The five scenarios used in the main experimental grid.}
    \label{tab:methodology-scenarios}
\end{table}

Each scenario's underlying profile is frozen ahead of time, and each carries a frozen set of energy and macronutrient targets computed independently of the pipeline under test, so every run's reported error is measured against the same target regardless of which variant or model produced the plan.

\section{Measurement controls in the evaluation harness}
\label{sec:methodology-harness}

No component of the harness scores its own work. The deterministic macro error is computed at collection time from the runtime's own hydration and totaller (\Cref{sec:design-eval-reuse}), written onto the run record, and never depends on a judge, so adding or rescoring a judge cannot move a quantitative number. Collection, scoring and reporting are separate stages joined only through stored artefacts (\Cref{fig:diagram-eval-flow}): the judge is handed a plan carrying no nutrient values and no indication of which variant or model produced it.

\begin{figure}[htbp]
    \centering
    \begin{tikzpicture}[
    stage/.style={draw, rounded corners, align=center, text width=50mm, minimum height=11mm, inner sep=2mm, font=\small},
    store/.style={draw, dashed, align=center, text width=46mm, minimum height=11mm, inner sep=1.6mm, font=\small\ttfamily},
    arr/.style={-{Latex[length=2mm]}},
    elab/.style={font=\small\itshape, text=black!70},
]

\node[stage] (s1) at (0, 0)
    {\textbf{run-cases}\\[0.6mm] drive Pipeline, compute macro error};
\node[store] (runs) at (0, -2.0) {runs/\{spec\_key\}.json};
\node[stage] (s2) at (0, -4.0)
    {\textbf{evaluate}\\[0.6mm] score each run, repeated per judge};

\node[store] (sc1) at (-3.8, -6.8)
    {scores\_judge\_1/\\[0.3mm]\{spec\_key\}\_\_jrep\{rep\}.json};
\node[store] (sc2) at ( 3.8, -6.8)
    {scores\_judge\_2/\\[0.3mm]\{spec\_key\}\_\_jrep\{rep\}.json};

\node[stage] (s3) at (0, -9.4)
    {\textbf{plot}\\[0.6mm] join artefacts and render figures};
\node[store] (figs) at (0, -11.4) {figures/, metrics/};

\draw[arr] (s1) -- node[elab, right] {generates} (runs);
\draw[arr] (runs) -- node[elab, right] {allows} (s2);
\draw[arr] (s2) -- node[elab, left, pos=0.45] {generates} (sc1.north);
\draw[arr] (s2) -- node[elab, right, pos=0.45] {generates} (sc2.north);
\draw[arr] (sc1) -- node[elab, left, pos=0.45] {allows} (s3);
\draw[arr] (sc2) -- node[elab, right, pos=0.45] {allows} (s3);

\draw[arr] (runs.east) -- node[elab, above] {allows} ++(40mm, 0) |- (s3.east);

\draw[arr] (s3) -- node[elab, right] {generates} (figs);

\end{tikzpicture}

    \caption{Data flow through the evaluation harness.}
    \label{fig:diagram-eval-flow}
\end{figure}

Two further controls keep a measured difference attributable to the flag that was varied. The same pipeline instance and food knowledge base connection are shared across every run of one model and variant, so no variant-specific initialisation difference -- a stale cache, a re-seeded index -- becomes a confound alongside that flag. A run's label is derived mechanically from the same configuration objects the runtime consumes (\Cref{sec:design-variant-config}), so it cannot drift from the configuration that produced it. Every reported number is regenerable from the stored per-run artefacts.

\section{Quantitative scoring: energy and macronutrient error}
\label{sec:methodology-quantitative}

The quantitative scoring stage measures how closely a plan's delivered energy and macronutrient totals match their targets, using the deterministic aggregator of \Cref{sec:design-totaller} to compute the delivered values -- independently of whether the totaller tool was available during generation, so that even a baseline-variant plan is judged by a real, exactly computed total.

For each tracked nutrient $n$ (energy, protein, carbohydrates, and fat) with a positive target, a signed relative error is computed as
\[
    \mathrm{rel}_n = \frac{\mathrm{actual}_n - \mathrm{target}_n}{\mathrm{target}_n} \times 100\%,
\]
where $\mathrm{actual}_n$ is the deterministically computed total for a hydrated plan and $\mathrm{target}_n$ is the profile's frozen target for that nutrient. Each run is then summarised by the mean absolute and the mean squared value of $\mathrm{rel}_n$.

\begin{definition}[Macro error]
    \label{def:macro-error}
    Let $N = \{\text{energy}, \text{protein}, \text{carbohydrate}, \text{fat}\}$ be the aggregated nutrients with a positive target for the run's profile. The \emph{macro error} of a run is the mean absolute relative error
    \[
        \mathrm{MAE\%} = \frac{1}{\left| N \right|} \sum_{n \in N} \left| \mathrm{rel}_n \right| .
    \]
    It is non-negative and unbounded above, and is zero exactly when every aggregated nutrient meets its target, so lower is better. It does not depend on a judge.
\end{definition}

\begin{definition}[Squared macro error]
    \label{def:macro-error-squared}
    With $N$ as in \Cref{def:macro-error}, the \emph{squared macro error} of a run is the mean squared relative error
    \[
        \mathrm{MSE\%} = \frac{1}{\left| N \right|} \sum_{n \in N} \mathrm{rel}_n^2 .
    \]
    It shares the properties above, but weights a single badly missed nutrient far more heavily than an even spread of small misses. It is always named explicitly where it is used, and likewise does not depend on a judge.
\end{definition}

Expressing the error relative to its target keeps a single figure comparable across nutrients on very different scales -- kilocalories against grams of fat. Both statistics are reported because this evaluation is most interested in the occasional badly wrong plan, not a uniform drift from target, and the squared form separates those cases.

\section{Qualitative scoring: the judge protocol}
\label{sec:methodology-qualitative}

Everything not captured by an energy or macronutrient number -- whether a recipe is physically coherent, whether a stated allergen was avoided, whether an explicit session request was honoured without violating a harder constraint -- is scored by an LLM-as-a-judge stage following the G-Eval pattern of \Cref{sec:background-evaluation}: a judge model is given an explicit natural-language rubric per criterion and asked to produce both a score and a short justification.

Every run is scored on one criterion present in all scenarios (whether the generated recipes are physically and logically coherent) plus every scenario-specific soft criterion listed for it in \Cref{sec:methodology-scenarios} -- for instance, whether a DASH-appropriate sodium-to-potassium balance was targeted, or whether a stated food preference was honoured without including a disliked food. Each criterion is scored on a continuous $[0, 1]$ scale, and the two judged metrics reported in \Cref{cha:results} are built from those scores as follows.

\begin{definition}[Soft score]
    \label{def:soft-score}
    Let $K$ be the set of soft criteria applying to a run: the recipe-coherence criterion present in every scenario, plus that scenario's own criteria (\Cref{sec:methodology-scenarios}). A judge scores each $k \in K$ with $s_k \in [0, 1]$, and the \emph{soft score} of one verdict is
    \[
        \mathrm{soft} = \frac{1}{\left| K \right|} \sum_{k \in K} s_k .
    \]
    The value carried into every figure and table is the mean of this over the three verdicts that judge returned for the run. It lies in $[0, 1]$, higher is better, and it is judge-dependent: it is reported once per judge and never averaged across judges.
\end{definition}

\begin{definition}[Safety score]
    \label{def:safety-score}
    The \emph{safety score} is a single always-scored rubric field in $[0, 1]$ recording whether the plan respected the profile's allergen exclusions and diet pattern, averaged over a judge's three verdicts exactly as the soft score is. It is deliberately not a member of $K$, so a serious safety failure cannot be diluted by the soft-score average. Each verdict also returns a list of the concrete \emph{safety violations} the judge identified, in free text. Higher is better, and it is judge-dependent.
\end{definition}

The judges are explicitly instructed not to penalise a plan for a missing micronutrient figure that the underlying food data itself does not cover, so that the coverage gaps described in \Cref{sec:impl-missing-nutrients} are not double-counted as a plan-quality failure on top of being flagged by the totaller's own coverage warnings.

\label{sec:impl-judges}
The initial scoring stage used a single judge and took its sampled score at face value, since the judge APIs expose no token probabilities from which G-Eval computes its probability-weighted score~\cite{liu2023geval}. Two problems surfaced. First, the judge is itself a sampled language model, so scoring the same plan twice does not reliably produce the same number. Second, the judge initially used belongs to the same model family as one of the compared configurations, exactly the situation in which the self-preference bias documented for this pattern~\cite{zheng2023judging} would inflate that configuration's scores. Neither is observable from inside a single-judge protocol. Both are addressed at the protocol level:

\begin{itemize}
    \item \emph{Two judges, from different vendors.} Every plan is scored independently by both. One shares a family with a compared configuration; the other has no such overlap. A family-specific bias therefore cannot hide, because it would appear as a divergence between two judges reading the same plans against the same rubric.
    \item \emph{Repeated scoring.} On the second of the two repetition axes, the verdicts a judge returns for a run are averaged before its qualitative score is used anywhere. Because the judge API does not expose token probabilities, this averaging also serves as a Monte Carlo estimate of G-Eval's probability-weighted score, which is the substitute G-Eval itself adopts whenever the scoring model's token probabilities are unavailable~\cite{liu2023geval}; it suppresses judge sampling noise without concealing how large it was.
    \item \emph{No averaging across judges.} Judge-dependent metrics -- the soft score and safety score -- are carried through the whole reporting pipeline as two parallel series and always reported side by side. Collapsing them into one number would destroy precisely the disagreement the second judge exists to reveal.
\end{itemize}

Each (run, judge, judge repetition) triple is persisted to its own file under a per-judge directory, so adding a judge re-scores nothing that already exists: the second judge was introduced after the first had scored the entire grid, and cost only its own calls. What the protocol revealed is reported in \Cref{sec:results-stability}.

\section{Stability metrics}
\label{sec:methodology-stability}

Repetition happens on two independent axes. Every cell is planned three times, sampling the generating model's nondeterminism; each run is then scored three times by each judge, sampling the judge's. The first axis is quantified here as run-to-run variability; the second is averaged away (\Cref{sec:methodology-qualitative}).

Two quantities describe how repeatable a result is, as opposed to how good it is, both reported per judge.

The first is \emph{run-to-run variability}: how much the same configuration's soft score moves when the identical request is planned again. Because the repetitions of one cell (\Cref{def:cell}) differ only in the model's sampling, the dispersion within a cell isolates that sampling noise.

\begin{definition}[Cell spread]
    \label{def:spread}
    The \emph{spread} of a cell is the mean absolute deviation of its soft score about the cell's own mean:
    \[
        \bar{s}_c = \frac{1}{R} \sum_{r=1}^{R} s_{c,r},
        \qquad
        d_c = \frac{1}{R} \sum_{r=1}^{R} \left| s_{c,r} - \bar{s}_c \right|,
    \]
    where:
    \begin{itemize}
        \item \( c \) -- a cell (\Cref{def:cell}),
        \item \( R \) -- the number of repetitions of a cell,
        \item \( s_{c,r} \) -- the soft score one judge awarded to repetition \( r \) of cell \( c \),
        \item \( \bar{s}_c \) -- the mean soft score of cell \( c \),
        \item \( d_c \) -- the spread of cell \( c \).
    \end{itemize}
    It is non-negative, expressed on the same $[0, 1]$ scale as the soft score, and lower means more repeatable. Like the soft score it is reported per judge.
\end{definition}

\begin{definition}[Run-to-run variability]
    \label{def:variability}
    The \emph{run-to-run variability} of a model, and of an experiment as a whole, is the mean cell spread over the relevant cells:
    \[
        V_m = \frac{1}{\left| C_m \right|} \sum_{c \in C_m} d_c,
        \qquad
        V = \frac{1}{\left| C \right|} \sum_{c \in C} d_c,
    \]
    where:
    \begin{itemize}
        \item \( C_m \) -- the set of cells whose language model is \( m \),
        \item \( C \) -- the set of all cells of the experiment,
        \item \( V_m \) -- the variability reported for model \( m \),
        \item \( V \) -- the pooled variability of the experiment.
    \end{itemize}
    Only cells with all of their repetitions scored contribute, so that no cell's spread is computed from a partial sample.
\end{definition}

The size of $V$ relative to the differences between configurations is what matters: a spread of the same order as those differences means they cannot be read as real.

The second is the agreement between the two judges. Because both score every plan under an identical rubric, any systematic difference between their means is a property of the judges, not the plans; any difference in how they \emph{rank} configurations is the more serious concern, since every comparative claim in \Cref{cha:results} depends on ranking. Both are reported in \Cref{sec:results-stability}.

\section{Operational metrics}
\label{sec:methodology-operational}

Four further figures are recorded per run without requiring any scoring step.

\begin{definition}[Iterations]
    \label{def:iterations}
    The \emph{iterations} of a run is the number of refinement passes the critique-and-refine loop completed. It is zero when reflection is disabled for the variant and also when the critic approved the first draft outright, and it is capped at the configured loop budget, three throughout this evaluation. It does not depend on a judge.
\end{definition}

\begin{definition}[Totaller calls]
    \label{def:totaller-calls}
    The \emph{totaller calls} of a run is the number of times the model invoked the deterministic totaller tool of \Cref{sec:design-totaller} while planning, counted across the nutrition agent and the refiner. It is zero when the totaller is disabled for the variant, and it counts only model-initiated invocations: the totals the critic is handed automatically (\Cref{sec:design-reflection}) and the aggregation the scoring stage performs itself (\Cref{sec:methodology-quantitative}) are not model decisions and are excluded. It is unbounded above and does not depend on a judge.
\end{definition}

\begin{definition}[Elapsed time]
    \label{def:elapsed-time}
    The \emph{elapsed time} of a run is its wall-clock duration in seconds, from the start of the planning pipeline to the delivered plan, including every model request and tool call the run made along the way. It does not depend on a judge.
\end{definition}

\begin{definition}[Token spend]
    \label{def:token-spend}
    The \emph{token spend} of a run is the number of tokens billed across every model request the run made, split into four additive categories matching the provider usage fields: cache read, which the provider served from a prompt cache instead of recomputing; input, the remainder of the billed input; output, the remainder of the billed output after reasoning tokens are set aside; and reasoning, the thinking tokens already included in the billed output. It does not depend on a judge and plays no role in scoring a plan's quality. No monetary cost is reported anywhere in this thesis: converting tokens into money requires per-model prices that change without notice and differ by routing provider, which would make any figure quoted here both unverifiable and out of date within months.
\end{definition}
